\documentclass[french,11pt]{article}

\usepackage{babel}
%\usepackage{fullpage}
\usepackage{psfig}
\usepackage{comment}
\usepackage{times}
\usepackage[T1]{fontenc}
\usepackage[french]{minitoc}
\usepackage[latin1]{inputenc}


\newenvironment{solution}{\nopagebreak\par\vspace*{2mm}\underline{Solution~:}\em\par}{\par\vspace*{2mm}}


\pagestyle{headings}
\markright{5/01/2006 \hfill Universit� Paris-Dauphine -- Paris \hfill\ } 

\baselineskip 6pt
\parskip 6pt
\onecolumn
\textwidth 16cm
\textheight 24cm
\hoffset=-2cm
\voffset=-2.5cm
\setlength{\parindent}{0.0cm}
\unitlength 1cm
\newtheorem{defin}{D�finition}[section]
\newtheorem{exampl}{Exemple}[section]
\def\picwidth{16cm}

\excludecomment{solution}

\begin{document}
\vspace*{0.5cm}

\begin{center}
\textbf{\Large Examen Bases de donn�es L3,  Jeudi 5 janvier 2006}\\
            Documents autoris�es
\end{center}

%%%%%%%%%%%%%%%%%%%%%%%%%%%%%%%%%%%%%%%%%%%%%%%%%%%%%%%%%%%%

\vspace*{0.5cm}

Un guide de voyage  d�cide de cr�er un site web proposant tous les restaurants
parisiens. Le site est construit sur une base de donn�es
contenant les restaurants, les plats propos�s et les menus.
Voici les tables.  Comme d'habitude, les cl�s �trang�res 
ne sont pas indiqu�es\,:
      \begin{itemize}
        \item {\tt Restaurant (\textbf{nom\_restau}, arrondissement, 
                    nom\_chef)}
         \item {\tt Plat (\textbf{id\_plat}, description, 
                          type\_plat)}
         \item {\tt Menu (\textbf{nom\_restau, no\_menu}, nom, prix)}
         \item {\tt D�tailMenu (\textbf{nom\_restau, no\_menu, id\_plat})} 
       \end{itemize}

L'attribut {\tt arrondissement} contient le num�ro de l'arrondissement du
restaurant.  La table {\tt Plat} donne l'ensemble des plats, tous
restaurants confondus. 
 Les types de plat sont (par exemple) {\tt Entr�e}, {\tt Plat
principal},  {\tt Dessert}, etc.  La table {\tt D�tailMenu} donne les plats
composant un menu\,: un ensemble d'entr�es, de plats principaux et de dessert.

\smallskip
{\bf Sch�ma (10 pts)}

\begin{enumerate}
  \item ({\bf 3 pts}) Vous d�cidez d'ouvrir � Dauphine  un restaurant
        nomm� ``Le Gratin''. Il propose 5 plats: carottes rap�es et
        huitres (entr�es), steak (plat principal), fruit et g�teau
        (dessert) et deux menus (``standard'' et ``gastro'') � 6 et 10 Euros.
        Le menu standard comprend  carottes rap�es, steak
                et  fruit. Le menu gastro comprend  huitres,
           steak et g�teau. Donnez le contenu des tables de 
           la base de donn�es pour repr�senter ce restaurant et ses menus.

   \item ({\bf 3 pts}) Donnez les commandes {\tt CREATE TABLE} 
         pour ce sch�ma. Indiquez 
         les cl\'es \'etrang\`eres et les cl\'es primaires.

   \item ({\bf 2 pts})
       Donner le sch�ma entit�/association correspondant � ce sch�ma.

   \item  ({\bf 2 pts}) On s'aper�oit qu'on a oubli� d'indiquer le
     choix de boissons propos� par un restaurant. Chaque boisson est
     caract�ris�e
        par son type (``vin blanc'', ``eau''), son nom (``Juran�on'',
        ``Vittel'')   et son prix. Proposez une extension du sch�ma
        E/A et
          de la base pour associer une liste de boissons � chaque restaurant.
\end{enumerate}

\smallskip
{\bf Requ�tes ({\bf 10 pts})}

La base contient tous les restaurants de Paris. Exprimez en SQL les requ�tes suivantes:

\begin{enumerate}
  \item  ({\bf 1 pt}) Nom des restaurants du douzi�me arrondissement.

  \item  ({\bf 1 pt}) Nom des restaurants proposant un plat
         dont la description contient le mot ``sauterelle''

  \item  ({\bf 2 pts})  Quels 
               chefs 
                 proposent des menus � 20 Euros (ne donner qu'une fois 
               chaque nom de chef)\,?
  
  \item  ({\bf 2 pts}) Donnez la liste des {\em entr�es} du menu 
   <<\,D�gustation\,>> du restaurant ``Le gratin''.

  \item  ({\bf 2 pts}) Quels restaurants (nom du restaurant, nom du chef) ne
    proposent pas de menu � moins de 20 Euros\,?

  \item  ({\bf 2 pts}) Quels menus du restaurant 
  ``Le gratin'' ont au moins un plat en commun (donner 
       les paires de noms de menus).

  \item  ({\bf 2 pts}) Donnez le nombre d'entr�es du menu
              ``Versailles'' du restaurant ``Boboche''.

  \end{enumerate}
  
\end{document}
